\section{Evaluation}
\label{sec:evaluation}

We propose a direct evaluation of the quality of \textsc{StyleDistance} embeddings, and an evaluation of their utility in downstream applications. We use other leading style representations like LISA \citep{lisa} and the embeddings from \citet{styleemb} as baselines, and compare them with our style embeddings on the STEL and STEL-or-Content tasks \citep{stel,styleemb}. We also compare to the LUAR model \citep{luar}, an authorship representation model which does not explicitly train content-independent representations but captures both content and style features in an attempt to represent different authors. LUAR is sometimes used for automatic style transfer evaluation, where we believe a strong style representation would be better suited. We thus choose this task for evaluation. % Finally, we evaluate whether our embeddings can be used to guide arbitrary text style transfer with the \textsc{TinyStyler} \citep{tinystyler} model.

\subsection{STEL and STEL-or-Content Evaluation}
\label{sec:steleval}


We first benchmark our embeddings on the STEL and STEL-or-Content tasks that allow for a direct evaluation of the quality of style representations. We briefly describe these tasks below and illustrate examples of these tasks in Appendix \ref{sec:appendix:stelfig}:

\begin{itemize}

\item \textbf{STEL:} Given two anchor sentences (A1, A2) and two test sentences (S1, S2), STEL measures the ability of an embedding model to pair each test sentence with the anchor sentence that shares the same style based on the cosine similarity of the embeddings of the sentences.

\item \textbf{STEL-or-Content:} The STEL-or-Content task is similar to STEL but more adversarially challenging, hence  better for testing content-independence. In this task, there are again two test sentences (S1, S2) but only a single anchor sentence (A). The test sentence that best matches the \emph{style} of the anchor must be selected; but the incorrect test sentence---which is written in a different style---is a  paraphrase of the anchor with similar \emph{content}. Therefore, in order to succeed on the STEL-or-Content task, a model needs to represent style features stronger than content features.

\end{itemize}



In their paper, \citet{stel} provide a STEL and STEL-or-Content evaluation benchmark over five features with curated natural data. We test our models on this benchmark and present the results in Table \ref{table:steleval}. Our results are consistent with results reported by \citet{stel}, who showed that even untrained models like \texttt{roberta-base} can capture style information well, resulting in stronger STEL performance than any of their fine-tuned models. However, the more challenging STEL-or-Content task, which better tests content-independence, shows that only models specifically trained for content-independence are able to capture style features better than content features. Our results indicate that the embeddings generated with our \textsc{StyleDistance} approach lead  over other style representations on both the STEL and STEL-or-Content tasks. Interestingly, we find that $\textsc{StyleDistance}_{\textsc{Synth}}$ captures style remarkably well, and manages to generalize to the natural text examples in the evaluation benchmark despite only being trained on our synthetic contrastive triplets. We conclude using synthetic parallel examples during training makes the model more content-independent and helps better capture style. 

\subsection{Generalization Experiment} 

In an ablation study, we test the ability of our training approach to produce embeddings that can generalize to unseen style features which are not present in the synthetic dataset. We conduct this evaluation by ablating features from the data used to train $\textsc{StyleDistance}_{\textsc{Synth}}$ under three conditions and show the results in Table \ref{table:ablationeval}. In the \textbf{In-Domain} condition, all 40 style features are included. In the \textbf{Out-of-Domain} condition, we exclude synthetic examples corresponding to the five style features in the STEL/STEL-or-Content benchmark \footnote{In our 40 features, there are two separate features for emoji and text emoticons (:-D) so we exclude 6 total features for this condition instead of 5 features, resulting in 34 features.}. In the \textbf{Out-of-Distribution} condition, we further exclude examples for any features similar or indirectly related to the five evaluated features. Details on the exact 15 style features ablated can be found in Appendix \ref{sec:appendix:ablationdetails}. We compare the Out-of-Domain and Out-of-Distribution performance of $\textsc{StyleDistance}_{\textsc{Synth}}$ on the two tasks to its In-Domain performance, obtained when it was trained on data for all 40 style features. Even in the challenging Out-of-Distribution condition, $\textsc{StyleDistance}_{\textsc{Synth}}$ retains 50\% of its performance on the challenging STEL-or-Content task (see the ``Retained Perf.'' column). This study indicates our training approach generalizes reasonably well to out-of-domain style features which can be composed from style features selected for generation and, to some extent, even to out-of-distribution style features fully outside the selected set.

\subsection{Synthetic Data for Probing}

Our previous experiments demonstrate that training on our synthetic dataset yields strong style embeddings. Next, we investigate whether the synthetic dataset can be used for an entirely different purpose: to probe which specific style features are captured by existing style representations. Our  \textsc{SynthStel} dataset allows for the creation of synthetic STEL and STEL-or-Content task instances across a  range of 40 style features---much broader than the \citet{stel} benchmark where five features were addressed. We use the test split of \textsc{SynthStel} to generate task instances for probing. We examine whether LISA vectors and the \citet{styleemb} style embeddings capture these 40 style features, which \textsc{StyleDistance} models are directly trained to represent. We show average results over all 40 features in Table \ref{table:synthsteleval}. Per feature results are provided in Appendix \ref{sec:appendix:fullsynthstelresults}. Our findings reveal only moderate coverage by LISA and \citet{styleemb}, with high variance depending on the evaluated feature (e.g., ``Usage of Nominalizations'' is poorly captured, with a near-zero STEL-or-Content score for both models). We calculate the mean squared error (MSE) between the STEL and STEL-or-Content scores for the real and synthetic task instances across the five features in the real benchmark, finding an average MSE of 0.039. This small MSE value shows that using synthetic data for probing can reasonably serve to assess which style features are represented by a model without need for manual example curation.

\synthstelevaltable

\subsection{Downstream Evaluation}

We next evaluate and/or demonstrate our style embeddings in three downstream applications.

\paragraph{Authorship Verification} We first test our style embeddings on the authorship verification (AV) task \citep{authorshipverification}. 
Given two documents by unknown authors, the goal of the authorship verification (AV) task is to determine whether they were written by the same author, based on their stylistic similarities and differences \citep{stylisticauthorshipverification}. We use a series of AV shared task datasets released by PAN in 2011-2015 \citep{pan11,pan13,pan14,pan15}\footnote{PAN is the ``Plagiarism Analysis, Authorship Identification, and Near-Duplicate Detection'' workshop. No AV shared task was proposed in 2012. (URL: \url{https://pan.webis.de})}. Since the two documents may be about different topics, good content-independent style representations would be expected to perform better in the AV task than embeddings that capture content. We calculate the cosine similarity of the two documents using our tested style embedding models with no fine-tuning, to measure their off-the-shelf ability to identify whether two documents were written by the same author and report results using the standard ROC-AUC metric used in AV. In Table \ref{table:avevaltable}, we compare the performance of \textsc{StyleDistance} embeddings against LISA and the \citet{styleemb} style embeddings. On average, \textsc{StyleDistance} outperforms the other representations on AV, demonstrating its effectiveness in representing style.

\avevaltable

\paragraph{Automatic Style Transfer Evaluation} \citet{lowresourceauthorshipstyletransfer} proposed the LUAR embedding model as an automatic measure for ``style transfer accuracy''. This approach was effective and was subsequently adopted for style transfer evaluation \citep{astropop,paraguide,tinystyler}. However, the LUAR model considers both style and content, hence confounding two aspects of style transfer evaluation---accuracy and meaning preservation---which are typically measured separately. Content-independent style representations would be a better measure for this task. We use the same evaluation dataset of 675 task instances used by  \citet{lowresourceauthorshipstyletransfer} where given an example text of a target author's style, the task is to discriminate which of two texts (a style transfer output and another actual text by the target author) is written by the target author. We show our results on this task in Table \ref{table:styletransfereval}. \textsc{StyleDistance} proves to be a more effective discriminator than all models, including LUAR. All models surpass human performance in distinguishing style transfer outputs. Since automatic style transfer evaluation typically includes a separate score for meaning preservation, using a model like LUAR (which is not content-independent) for measuring style transfer accuracy undermines the rigor of the style transfer accuracy metric. We find that a robust content-independent model like \textsc{StyleDistance} may enhance automatic style transfer evaluation by: (1) acting as a stronger discriminator, and (2) ensuring style transfer accuracy is assessed independently of meaning preservation.

\paragraph{Style Transfer Steering} Previous systems, like TinyStyler, have leveraged style embeddings to steer style transfer \citep{tinystyler}. While the original TinyStyler system rewrites text by conditioning on \citet{styleemb} embeddings, we demonstrate that \textsc{StyleDistance} provides an alternative, and reproduce TinyStyler with \textsc{StyleDistance} embeddings. We showcase an example of an output in Table \ref{table:smalltinystyleroutputs} with more details and results in Appendix \ref{sec:appendix:tinystyler}. We will make this version of TinyStyler available as a resource. With this result, we demonstrate \textsc{StyleDistance} can be used as a simple drop-in replacement for downstream applications in systems where weaker style representations have been previously used.
\styletransferevaltable
\smalltinystyleroutputs
